\documentclass[11pt]{letter}
\usepackage[utf8]{inputenc}
\usepackage[T1]{fontenc}
\usepackage{lmodern}
\usepackage[margin=2.5cm]{geometry}
\signature{Richar Andre Vilca-Solorzano \\ (on behalf of all authors)}
\address{Faculty of Statistical and Computer Engineering\\
Universidad Nacional del Altiplano (UNAP)\\ Puno, Peru\\
\texttt{75521963@est.unap.edu.pe}}
\begin{document}
\begin{letter}{The Editor-in-Chief\\ \textit{Remote Sensing}\\ MDPI}
\opening{Dear Editor,}

We are pleased to submit our manuscript, \textit{``Harmonized Sentinel-2/Landsat
Retrieval of Water Transparency in Lake Titicaca Using Machine Learning: A
Transferable Validation and Retrievability Protocol for High-Altitude Lakes
(2013--2024)''}, for consideration as a research article in \textit{Remote Sensing}.

The study sits squarely within the journal's scope---machine-learning retrieval of
water-quality variables from harmonized optical satellite missions. We present, to
our knowledge, the first multi-mission harmonized retrieval of water transparency
(Secchi depth) for a high-altitude Andean lake (3{,}810~m a.s.l.), built entirely
on \textbf{real in-situ measurements} (881 field observations, 1{,}002
satellite--field match-ups, 2013--2024) sampled at their true coordinates across
the whole lake. It directly complements recent work on Lake Titicaca published in
\textit{Remote Sensing} using global products (Maligaya et al., 2024), to which we
add harmonized multi-mission retrieval, local in-situ validation, and an explicit
retrievability assessment.

We believe this work makes four distinctive contributions. (1) \textbf{A
transferable validation hierarchy:} beyond station-wise cross-validation and a
temporal hold-out, we introduce a leave-one-zone-out test showing that, although
random and station-wise validation are identical (no spatial-autocorrelation
leakage), extrapolation to an unseen optical regime collapses $R^2$ toward zero---a
loss of transferability that random splits conceal. (2) \textbf{Calibrated
uncertainty:} we attach per-pixel split-conformal prediction intervals (empirical
coverage 89.8\% at a 90\% nominal level). (3) \textbf{Physical interpretability:}
SHAP shows the retrieval is driven by green/blue reflectance ratios, consistent
with water-clarity optics. (4) \textbf{Transparent delimitation of limits:} we
report explicitly that chlorophyll-$a$, suspended solids and thermal temperature
are \emph{not} retrievable in this oligotrophic regime---an honest operational
envelope that cautions against transferring turbid-water algorithms to clear
high-altitude lakes. A lake-wide transparency map demonstrates operational utility.

All data and code are openly available, and the full pipeline is reproducible from
public satellite archives and published institutional monitoring records. The
manuscript is original, not under consideration elsewhere, and all authors approve
the submission. We have no conflicts of interest to declare.

Thank you for considering our work.

\closing{Sincerely,}
\end{letter}
\end{document}
